\documentclass[12pt,a4paper]{article}

\usepackage{iftex}

\ifPDFTeX
  \usepackage{cmap}
  \usepackage[T2A]{fontenc}
  \usepackage[utf8]{inputenc}
  \usepackage{lmodern}
  \usepackage[english,russian]{babel}
\else
  \usepackage{fontspec}
  \defaultfontfeatures{Ligatures=TeX}
  \setmainfont{Times New Roman}
  \setsansfont{Arial}
  \setmonofont{Consolas}
  \usepackage{polyglossia}
  \setmainlanguage{russian}
  \setotherlanguage{english}
\fi

\usepackage{amsmath,amssymb,amsthm,mathtools}
\usepackage[a4paper,margin=2.2cm,top=2.4cm,bottom=2.2cm]{geometry}
\usepackage{enumitem}
\usepackage{fancyhdr}
\usepackage[hidelinks]{hyperref}

\providecommand{\SheetNo}{}
\newcommand{\sheetauthor}{}

\newcommand{\sheettitle}[1]{
  \begin{center}
    {\Large\bfseries Серия \SheetNo. #1}\\[0.4em]
    \rule{\linewidth}{0.4pt}
  \end{center}
  \vspace{0.5em}
  \markboth{#1}{#1}
}

\setlength{\headheight}{15pt}
\pagestyle{fancy}
\fancyhf{}
\fancyhead[L]{\small Серия \SheetNo}
\fancyhead[R]{\small\leftmark}
\fancyfoot[C]{\thepage}
\renewcommand{\headrulewidth}{0.4pt}

\newlist{problems}{enumerate}{2}
\setlist[problems]{label=\arabic*., ref=\arabic*,
                   leftmargin=*, itemsep=0.8em, topsep=0.8em}

\newcommand{\cyralphnum}[1]{
  \ifcase#1\or а\or б\or в\or г\or д\or е\or ж\or з\or и\or к\or л\or м\or
  н\or п\or р\or с\or т\or у\or ф\or х\or ц\or ч\or ш\or щ\or э\or ю\or я\fi}
\newcommand{\cyralph}[1]{\cyralphnum{\value{#1}}}
\AddEnumerateCounter{\cyralph}{\cyralphnum}{щ}

\newlist{parts}{enumerate}{2}
\setlist[parts]{label=\cyralph*), ref=\cyralph*,
                leftmargin=1.5em, itemsep=0.25em, topsep=0.25em}

\newcommand{\N}{\mathbb{N}}
\newcommand{\Z}{\mathbb{Z}}
\newcommand{\Q}{\mathbb{Q}}
\newcommand{\R}{\mathbb{R}}
\providecommand{\C}{}
\renewcommand{\C}{\mathbb{C}}
\providecommand{\NOD}{}
\renewcommand{\NOD}{\operatorname{\text{НОД}}}
\providecommand{\NOK}{}
\renewcommand{\NOK}{\operatorname{\text{НОК}}}


\def\SheetNo{1}

\begin{document}

\sheettitle{Введение. Учимся доказывать.}

\begin{problems}

\item Докажите, что
\[
1\cdot 1!+2\cdot 2!+\ldots+n\cdot n!=(n+1)!-1.
\]

\item Докажите формулу
\[
1^2+2^2+\ldots+n^2=\frac{n(n+1)(2n+1)}{6}.
\]

\item Найдите ошибку в следующем доказательстве утверждения «все лошади одного цвета». База индукции при $n=1$ очевидна. Пусть в любом множестве из $n$ лошадей все лошади одного цвета. Рассмотрим множество из $n+1$ лошадей $h_1,h_2,\ldots,h_{n+1}$. По предположению индукции лошади $h_1,\ldots,h_n$ одного цвета, и лошади $h_2,\ldots,h_{n+1}$ тоже одного цвета. Так как эти два множества пересекаются, все $n+1$ лошадей одного цвета.

\item Последовательность задана условиями $u_1=2$, $u_2=3$ и $u_n=3u_{n-1}-2u_{n-2}$ при $n>2$. Докажите, что $u_n=2^{n-1}+1$.

\item Докажите для всех $n\in\N$:
\begin{parts}
  \item $6\mid(n^3-n)$;
  \item $30\mid(n^5-n)$.
\end{parts}

\item Найдите и докажите формулу для суммы
\[
1\cdot 2^1+2\cdot 2^2+\ldots+n\cdot 2^n.
\]

\item Докажите, что при любом $n\in\N$:
\begin{parts}
  \item
  \[
  1+\frac12+\frac1{2^2}+\ldots+\frac1{2^n}<2;
  \]
  \item
  \[
  \prod_{k=1}^{n}\left(1+\frac1{2^k}\right)<\frac52.
  \]
\end{parts}

\item Пусть $x\ne0$ и число $x+x^{-1}$ целое. Докажите, что $x^n+x^{-n}$ целое при любом $n\in\N$.

\item Последовательность Фибоначчи задана условиями $F_0=0$, $F_1=1$, $F_{n+1}=F_n+F_{n-1}$. Докажите для $m,n\geqslant1$, что
\[
F_{m+n}=F_{m-1}F_n+F_mF_{n+1},
\]
и получите из этого тождество Кассини.

\item Докажите, что для любого $n\geqslant3$ единицу можно представить в виде суммы $n$ различных дробей с числителем 1 и натуральными знаменателями.

\item Назовём натуральное число ровным, если все цифры его десятичной записи одинаковы. Докажите, что любое $n$-значное натуральное число представимо в виде суммы не более чем $n+1$ ровных чисел.

\item На доске записана строка длины 100, состоящая из нулей и единиц. Разрешается либо заменить первую цифру, либо заменить цифру, стоящую сразу после первой единицы. Докажите, что из любой булевой строки можно получить любую другую булевую строку.

\item Выпуклый $n$-угольник разбит непересекающимися диагоналями на треугольники. Докажите, что в разбиении ровно $n-2$ треугольника и $n-3$ диагонали, независимо от способа разбиения.

\item На плоскости проведено несколько прямых общего положения (никакие две из них не параллельны и никакие три не пересекаются в одной точке).
\begin{parts}
  \item На сколько частей они разбивают плоскость?
  \item Может ли среди частей, на которые они разбивают плоскость, оказаться ровно 239 многоугольников?
\end{parts}

\item На фестиваль военно-морской песни приглашены хоры из 100 стран. Каждый хор должен исполнить три песни, после чего сразу уехать домой (не выслушивая следующие песни). Ознакомившись с текстами песен, организаторы обнаружили, что каждая песня оскорбительна для одной из участвующих стран. Докажите, что порядок выступлений можно назначить так, чтобы никто не выслушивал более трех оскорбительных для его страны песен.

\item На листе клетчатой бумаги с размером клетки $1\times1$ изображен прямоугольник. Прямоугольник разбит прямыми, параллельными его сторонам, на некоторое количество маленьких прямоугольников. У каждого маленького прямоугольника длины сторон выражаются целыми числами, при этом длина хотя бы одной его стороны чётна. Докажите, что длина хотя бы одной стороны исходного прямоугольника также является чётным числом.

\item В каждой клетке таблицы $10\times10$ стоит одна из цифр. Изначально во всех клетках таблицы стоят нули. У каждой строки и у каждого столбца есть кнопка, которая увеличивает все числа в этой строке или столбце на 1; при этом девятки заменяются на нули. Петя понажимал на какие-то кнопки. Докажите, что таблицу можно вернуть к первоначальному состоянию, нажав на кнопки:
\begin{parts}
  \item не более 180 раз;
  \item не более 90 раз.
\end{parts}

\item Взяли несколько одинаковых равносторонних треугольников. Вершины каждого из них пометили цифрами 1, 2 и 3. Затем их сложили в стопку. Могло ли оказаться, что сумма чисел, находящихся в каждом углу, равна 55?

\item В олимпиаде участвовали $n$ человек. Все решали одинаковые $k$ задач. Будем называть задачу простой, если её решили больше половины человек. Может ли оказаться, что все задачи были простыми, но каждый человек решил меньше половины задач?

\item В цирке выступает $n$ клоунов. Каждый клоун оделся, использовав в элементах одежды не менее 5 из имеющихся 12 цветов. Известно, что никакие два клоуна не оделись в одинаковый набор цветов, а также никакой цвет не использован более 20 раз. Найдите наибольшее возможное значение $n$.

\item В классе 20 школьников. Было устроено несколько экскурсий, в каждой из которых участвовало хотя бы четверо школьников этого класса. Докажите, что найдётся такая экскурсия, что каждый из участвовавших в ней школьников принял участие по меньшей мере в $1/17$ всех экскурсий.

\item Докажите, что в любой компании из 10 человек найдутся либо трое попарно знакомых, либо четверо попарно незнакомых.

\item В стае обезьян 80 детёнышей. У любых двух из них есть общий дед и, разумеется, у каждого детёныша два деда. Докажите, что найдётся дед, у которого хотя бы 54 внука.

\item Даны 10 двузначных чисел. Докажите, что среди данных чисел можно выбрать два непустых непересекающихся множества так, чтобы суммы чисел в них были равны. Некоторые числа можно вообще не использовать.

\item В ряд выписаны натуральные числа $a_1,a_2,\ldots,a_7$. Докажите, что можно выбрать некоторые из них и поставить перед каждым выбранным числом $+$ или $-$ так, чтобы значение полученного выражения делилось на 101.

\item Каждый из учеников класса занимается не более чем в двух кружках, причём для любой пары учеников существует кружок, в котором они занимаются вместе. Докажите, что найдётся кружок, в котором занимается не менее $2/3$ всего класса.

\item В классе 28 учеников. На уроке программирования они делятся на три группы. На уроке английского языка они тоже делятся на три группы, но по-другому. И на уроке физкультуры они делятся на три группы каким-то третьим способом. Докажите, что найдутся хотя бы два ученика, которые на всех трёх занятиях находятся друг с другом в одной группе.

\item По краю круглого стола равномерно расставлены таблички с фамилиями дипломатов, участвующих в переговорах. После начала переговоров оказалось, что ни один из дипломатов не сидит против своей таблички. Можно ли повернуть стол так, чтобы по крайней мере два дипломата сидели против своих табличек?

\item Каждое натуральное число покрасили в один из трёх цветов: красный, синий или зелёный, причём все три цвета встречаются. Может ли оказаться так, что сумма любых двух чисел разных цветов является числом оставшегося цвета?

\item На столе лежат 15 салфеток произвольной формы и размера, которые его полностью закрывают. Докажите, что можно убрать 8 из них так, чтобы оставшиеся закрывали не менее $7/15$ площади стола.

\end{problems}

\end{document}
